\documentclass[11pt]{article}
\usepackage[T1]{fontenc}
\usepackage{lmodern}
\usepackage{microtype}
\usepackage{amsmath, amssymb}
\usepackage{geometry}
\geometry{margin=1in}
\usepackage{hyperref}
\hypersetup{colorlinks=true, linkcolor=blue, urlcolor=blue, citecolor=blue}

\title{Hypergraph Placement Model for R/A/B/C Hamiltonian Application}
\author{Uni20 Developer Documentation}
\date{\today}

\begin{document}

\maketitle

\section{Scope}

This note defines the scheduling domain and cost functional for applying a
block-sparse effective Hamiltonian in the R/A/B/C form used by the current
TensorContraction bridge.  The immediate target is the two-site DMRG Lanczos
operator
\begin{equation}
  R_r \mathrel{+}= f_t A_a B_b C_c,
  \qquad t = (r,a,b,c,f_t),
\end{equation}
where $B$ is the input Krylov vector block family and $R$ is the output Krylov
vector block family.

For Hamiltonian application inside Lanczos, $B$ and $R$ must have the same
logical block space and the same persistent device layout.  The next Lanczos
matvec consumes the previous output as its input, so an implicit relayout after
each application is not acceptable.  The same R/A/B/C machinery can also be
used for environment construction, but then the input and output families are
not necessarily the same object and need not share a device layout.  This
document records that distinction, but focuses first on the Hamiltonian case.

\section{The Sparse Coefficient Tensor as a Hypergraph}

Let
\begin{itemize}
  \item $\mathcal{D}$ be the set of CUDA devices visible to one MPI rank,
  \item $\mathcal{K}$ be the logical center-block set shared by $B$ and $R$,
  \item $\mathcal{A}$ and $\mathcal{C}$ be the left and right environment block sets,
  \item $\mathcal{T}$ be the set of nonzero R/A/B/C terms.
\end{itemize}
Each term
\begin{equation}
  t \in \mathcal{T}, \qquad t = (r(t),a(t),b(t),c(t),f_t)
\end{equation}
is a typed hyperedge connecting one output block, one input center block, and
one left/right environment pair.  Symmetry explains much of the sparsity, but
the optimizer should use the actual nonzero tensor $f$ as the ground-truth
connectivity object.

\section{Hamiltonian Scheduling Domain}

The first useful Hamiltonian domain consists of four classes of choices:
center-block ownership, active environment placement, exact factorization of
the sparse coefficient tensor into temporary families, and placement of the
first-stage product contributions and second-stage operations.

\paragraph{Center-block owner.}
Choose one canonical owner for each center block:
\begin{equation}
  h : \mathcal{K} \to \mathcal{D}.
\end{equation}
For Hamiltonian application, $h(k)$ is both the persistent owner of input block
$B_k$ and the final owner of output block $R_k$.

\paragraph{Active environment placement.}
Choose where active environment blocks are resident during the local solve:
\begin{align}
  e_A &: \mathcal{A} \to \mathcal{D}, \\
  e_C &: \mathcal{C} \to \mathcal{D}.
\end{align}
The first model treats each active environment block as having one selected
device.  A later model may allow environment replicas, but then the second-stage
device must choose one of the available replicas.

\paragraph{Exact path factorization.}
Split the sparse coefficient tensor into a left-first part and a right-first
part:
\begin{equation}
  f_{rabc} = f^L_{rabc} + f^R_{rabc}.
\end{equation}
Each part is then represented by an exact sparse factorization through
temporary labels:
\begin{align}
  f^L_{rabc} &= \sum_{p \in \mathcal{P}_L}
    \lambda^L_{pab}\,\mu^L_{rpc}, \\
  f^R_{rabc} &= \sum_{q \in \mathcal{P}_R}
    \lambda^R_{qbc}\,\mu^R_{raq}.
\end{align}
This induces left-first temporaries
\begin{equation}
  P^L_p = \sum_{a,b} \lambda^L_{pab} A_a B_b
\end{equation}
and right-first temporaries
\begin{equation}
  P^R_q = \sum_{b,c} \lambda^R_{qbc} B_b C_c.
\end{equation}
The final output is
\begin{equation}
  R_r \mathrel{+}=
    \sum_{p,c} \mu^L_{rpc} P^L_p C_c
  + \sum_{a,q} \mu^R_{raq} A_a P^R_q .
\end{equation}
All products accumulated into one temporary must have the same dense matrix
shape, and all downstream uses must be shape-compatible with the target
output block.  The factorization is therefore typed by both block labels and
matrix dimensions; it is not merely an algebraic factorization of scalar
coefficients.

This formulation incorporates the scalar prefactor before temporary
accumulation.  For example, if two first-stage products are reused with
opposite signs downstream, the exact factorization may create
\begin{equation}
  P^L_1 = A_1B_1 + A_2B_2,
  \qquad
  P^L_2 = A_1B_1 - A_2B_2,
\end{equation}
so that different $C$ blocks can multiply different linear combinations.  This
is a symbolic sparse factorization, not a numerical low-rank approximation.
The first optimizer should not use SVD, QR, or any other floating-point rank
reduction of $f$ to create the $P$ objects: even if such a reduction lowers the
formal rank, it changes numerical stability and reproducibility properties.

The term-level model is a degenerate exact factorization.  A right-first term
$(r,a,b,c,f_t)$ can be represented with one temporary $q=(b,c)$ and coefficient
$\mu^R_{raq}=f_t$, or with a private temporary $q=t$.  The useful
factorization groups terms only when doing so creates an exactly reusable
linear combination.

\paragraph{Conservative no-rank-reduction specialization.}
Because we do not currently know a guaranteed stable useful factorization of
$f$, the first implementation should use a conservative exact specialization.
For left-first terms, use
\begin{equation}
  p=(r,c), \qquad
  \lambda^L_{(r,c)ab} = f^L_{rabc}, \qquad
  \mu^L_{r'p c'} =
    \begin{cases}
      1, & p=(r',c'),\\
      0, & \text{otherwise}.
    \end{cases}
\end{equation}
Thus
\begin{equation}
  P^L_{rc} = \sum_{a,b} f^L_{rabc} A_aB_b,
  \qquad
  R_r \mathrel{+}= P^L_{rc}C_c.
\end{equation}
For right-first terms, use
\begin{equation}
  q=(r,a), \qquad
  \lambda^R_{(r,a)bc} = f^R_{rabc}, \qquad
  \mu^R_{r'a'q} =
    \begin{cases}
      1, & q=(r',a'),\\
      0, & \text{otherwise}.
    \end{cases}
\end{equation}
Thus
\begin{equation}
  P^R_{ra} = \sum_{b,c} f^R_{rabc} B_bC_c,
  \qquad
  R_r \mathrel{+}= A_aP^R_{ra}.
\end{equation}
This specialization preserves the general $\lambda/\mu$ notation while avoiding
any numerical factorization.  It also fixes the natural second-stage
destination: $P^L_{rc}$ must be assembled where $C_c$ resides, and $P^R_{ra}$
must be assembled where $A_a$ resides.

\paragraph{Product and second-stage placement.}
Choose the device for each nonzero first-stage product contribution:
\begin{align}
  \kappa_L(p,a,b) &\in \mathcal{D}
    &&\text{when } \lambda^L_{pab} \ne 0,\\
  \kappa_R(q,b,c) &\in \mathcal{D}
    &&\text{when } \lambda^R_{qbc} \ne 0.
\end{align}
These choices create device-labelled partial temporaries.  Also choose the
device for each nonzero second-stage use:
\begin{align}
  \sigma_L(r,p,c) &\in \mathcal{D}
    &&\text{when } \mu^L_{rpc} \ne 0,\\
  \sigma_R(r,a,q) &\in \mathcal{D}
    &&\text{when } \mu^R_{raq} \ne 0.
\end{align}
For the conservative specialization above, the first natural restriction is
\begin{align}
  \sigma_L(r,(r,c),c) &= e_C(c),\\
  \sigma_R(r,a,(r,a)) &= e_A(a).
\end{align}
If a product contribution is accumulated on a different device from a
second-stage consumer of its temporary, that partial temporary must be migrated
and accumulated on the consumer device.  Thus temporary migration is not an
optional extension: it is part of the induced execution DAG.

The tuple
$(h,e_A,e_C,\mathcal{P}_L,\mathcal{P}_R,\lambda,\mu,\kappa,\sigma)$ is still
not the fully general CUDA scheduling space.  It excludes arbitrary temporary
reduction trees and nontrivial final-output reduction trees.  These exclusions
keep the first optimizer tractable while capturing the main placement
tradeoffs.

\section{Induced Execution DAG}

Given the placement and factorization tuple above, the execution DAG is
determined by the following canonical rules.

\paragraph{Input-center availability.}
For each center block $b$, define the devices that need $B_b$:
\begin{equation}
  U_B(b) =
    \{ \kappa_L(p,a,b) : \lambda^L_{pab}\ne 0 \text{ for some } p,a \}
    \cup
    \{ \kappa_R(q,b,c) : \lambda^R_{qbc}\ne 0 \text{ for some } q,c \}.
\end{equation}
The block is already resident on $h(b)$.  If $d \in U_B(b)$ and $d \ne h(b)$,
the model includes a peer copy or staging operation for $B_b$ to device $d$.
There is no reason to replicate $B_b$ onto a device that has no assigned term
using it.

\paragraph{Environment availability.}
Similarly, define
\begin{align}
  U_A(a) &=
    \{ \kappa_L(p,a,b) : \lambda^L_{pab}\ne 0 \text{ for some } p,b \}
    \cup
    \{ \sigma_R(r,a,q) : \mu^R_{raq}\ne 0 \text{ for some } r,q \}, \\
  U_C(c) &=
    \{ \kappa_R(q,b,c) : \lambda^R_{qbc}\ne 0 \text{ for some } q,b \}
    \cup
    \{ \sigma_L(r,p,c) : \mu^L_{rpc}\ne 0 \text{ for some } r,p \}.
\end{align}
For steady-state resident Lanczos modelling, active $A$ and $C$ blocks may be
assumed already materialized on the devices that need them.  For cold-start or
environment-staging models, the cost functional should include the required
$A$ and $C$ staging bytes.

\paragraph{First-stage temporaries.}
Each device constructs local partial temporaries from the first-stage products
assigned to that device:
\begin{align}
  P^{L,d}_p &=
    \sum_{\substack{a,b\\ \lambda^L_{pab}\ne 0\\ \kappa_L(p,a,b)=d}}
      \lambda^L_{pab} A_aB_b,\\
  P^{R,d}_q &=
    \sum_{\substack{b,c\\ \lambda^R_{qbc}\ne 0\\ \kappa_R(q,b,c)=d}}
      \lambda^R_{qbc} B_bC_c.
\end{align}
The coefficients $\lambda$ are applied during this accumulation, so identical
matrix products that appear with different signs or scalars can feed distinct
temporary objects.

\paragraph{Temporary migration.}
For each left-first temporary, define its consumer devices
\begin{equation}
  V_L(p) = \{ \sigma_L(r,p,c) : \mu^L_{rpc}\ne 0 \}.
\end{equation}
For each right-first temporary, define
\begin{equation}
  V_R(q) = \{ \sigma_R(r,a,q) : \mu^R_{raq}\ne 0 \}.
\end{equation}
If a nonzero partial $P^{L,d}_p$ exists and $v \in V_L(p)$ with $v\ne d$, then
$P^{L,d}_p$ is copied to $v$ and accumulated into the assembled temporary on
$v$.  The analogous rule applies to $P^{R,d}_q$.  A later scheduler may choose
tree or multicast strategies, but the first model counts one materialized
partial-temporary transfer per source and consumer device.

\paragraph{Second-stage accumulation.}
Each nonzero $\mu^L_{rpc}$ computes and accumulates
$\mu^L_{rpc}P^{L,\sigma_L(r,p,c)}_pC_c$ on $\sigma_L(r,p,c)$.  Each nonzero
$\mu^R_{raq}$ computes and accumulates
$\mu^R_{raq}A_aP^{R,\sigma_R(r,a,q)}_q$ on $\sigma_R(r,a,q)$.  The result is
accumulated into a per-device partial output block
\begin{equation}
  P^d_r.
\end{equation}

\paragraph{Final output reduction.}
For each output block $r$, all nonzero partials $P^d_r$ are reduced or copied
to the canonical owner $h(r)$.  If $d=h(r)$, the partial can be accumulated
directly into the final $R_r$ buffer.  If $d \ne h(r)$, the model includes a
peer transfer and accumulation into $R_r$ on $h(r)$.

\section{Cost Functional}

The objective is predicted makespan.  For each device $d$, define a device cost
\begin{equation}
  C_d =
    C^{\mathrm{gemm}}_d
  + C^{\mathrm{copy}}_d
  + C^{\mathrm{overhead}}_d
  + C^{\mathrm{memory}}_d.
\end{equation}
The optimizer minimizes
\begin{equation}
  \min \max_{d \in \mathcal{D}} C_d,
\end{equation}
possibly subject to memory-capacity constraints.

\subsection{GEMM Work}

Let $\Phi^{L,1}_{ab}$ be the flop count for $A_aB_b$, and
$\Phi^{R,1}_{bc}$ be the flop count for $B_bC_c$.  Let
$\Phi^{L,2}_{pc}$ be the flop count for $P^L_pC_c$, and let
$\Phi^{R,2}_{aq}$ be the flop count for $A_aP^R_q$.  Then
\begin{align}
  C^{\mathrm{gemm}}_d &=
    \alpha_{L,1}
      \sum_{\substack{p,a,b\\ \lambda^L_{pab}\ne 0\\ \kappa_L(p,a,b)=d}}
        \Phi^{L,1}_{ab}
  + \alpha_{R,1}
      \sum_{\substack{q,b,c\\ \lambda^R_{qbc}\ne 0\\ \kappa_R(q,b,c)=d}}
        \Phi^{R,1}_{bc} \notag \\
  &\quad
  + \alpha_{L,2}
      \sum_{\substack{r,p,c\\ \mu^L_{rpc}\ne 0\\ \sigma_L(r,p,c)=d}}
        \Phi^{L,2}_{pc}
  + \alpha_{R,2}
      \sum_{\substack{r,a,q\\ \mu^R_{raq}\ne 0\\ \sigma_R(r,a,q)=d}}
        \Phi^{R,2}_{aq}.
\end{align}
The coefficients $\alpha$ are not just reciprocal peak FLOP rates.  For small
block GEMMs they should absorb cuBLAS dispatch overhead and the reduced
efficiency of small matrices.

Temporary construction also includes accumulation operations.  For large dense
blocks these are lower order than GEMMs, but for many tiny symmetry blocks they
can contribute measurable kernel-launch and memory-bandwidth overhead.  The
overhead model below should therefore count the number of product accumulations
into each temporary, not only the number of final temporary objects.

\subsection{Data Movement}

Let $\operatorname{bytes}(M_i)$ be the storage size of block $M_i$.  The
steady-state input-center movement cost is
\begin{equation}
  C^{B\mathrm{-copy}}_d =
  \beta_B
  \sum_{\substack{b \in \mathcal{K}\\d \in U_B(b)\\d \ne h(b)}}
    \operatorname{bytes}(B_b).
\end{equation}
Cold-start environment staging can be modelled as
\begin{equation}
  C^{AC\mathrm{-stage}}_d =
  \beta_A \sum_{\substack{a \in \mathcal{A}\\d \in U_A(a)}} \operatorname{bytes}(A_a)
  +
  \beta_C \sum_{\substack{c \in \mathcal{C}\\d \in U_C(c)}} \operatorname{bytes}(C_c).
\end{equation}
This term should normally be omitted or amortized when modelling repeated
resident Lanczos matvecs with fixed active environments.
If active environments are already resident according to $e_A$ and $e_C$, then
only additional replicas should be charged:
\begin{equation}
  C^{AC\mathrm{-replica}}_d =
  \beta_A \sum_{\substack{a \in \mathcal{A}\\d \in U_A(a)\\d\ne e_A(a)}} \operatorname{bytes}(A_a)
  +
  \beta_C \sum_{\substack{c \in \mathcal{C}\\d \in U_C(c)\\d\ne e_C(c)}} \operatorname{bytes}(C_c).
\end{equation}
Which of these two terms is appropriate depends on whether the model is
scoring cold environment materialization or steady-state resident matvecs.

The final output movement cost is
\begin{equation}
  C^{R\mathrm{-reduce}}_d =
  \rho_R
  \sum_{\substack{r \in \mathcal{K}\\P^d_r\ne 0\\d \ne h(r)}}
    \operatorname{bytes}(R_r).
\end{equation}
In the restricted domain, this is a flat peer-copy-plus-accumulate model.  A
future model may replace it with an explicit reduction tree.

Temporary migration is
\begin{align}
  S_L(p) &=
    \{ \kappa_L(p,a,b) : \lambda^L_{pab}\ne 0 \},\\
  S_R(q) &=
    \{ \kappa_R(q,b,c) : \lambda^R_{qbc}\ne 0 \}.
\end{align}
The corresponding movement cost is
\begin{align}
  C^{P\mathrm{-copy}}_d &=
  \beta_P
  \left(
    \sum_{\substack{p\in \mathcal{P}_L\\s\in S_L(p)\\d\in V_L(p)\\d\ne s}}
      \operatorname{bytes}(P^L_p)
    +
    \sum_{\substack{q\in \mathcal{P}_R\\s\in S_R(q)\\d\in V_R(q)\\d\ne s}}
      \operatorname{bytes}(P^R_q)
  \right).
\end{align}
This term is essential.  A schedule that saves first-stage GEMMs but ships a
large temporary to many devices may be worse than recomputing a smaller
temporary locally.

\subsection{CUDA and Scheduling Overheads}

CUDA overheads are not represented well by FLOPs or bytes alone.  The cost
model should include count features such as
\begin{equation}
  C^{\mathrm{overhead}}_d =
    \gamma_T N^{\mathrm{term}}_d
  + \gamma_L N^{L\mathrm{-prod}}_d
  + \gamma_R N^{R\mathrm{-prod}}_d
  + \gamma_{PC} N^{P\mathrm{-copy}}_d
  + \gamma_{PR} N^{\mathrm{partial}}_d
  + \gamma_Q N^{\mathrm{peer}}_d,
\end{equation}
where $N^{L\mathrm{-prod}}_d$ and $N^{R\mathrm{-prod}}_d$ count first-stage
product accumulations on $d$, $N^{P\mathrm{-copy}}_d$ counts temporary
migrations, $N^{\mathrm{partial}}_d$ counts nonzero partial output blocks on
$d$, and $N^{\mathrm{peer}}_d$ counts peer-staged input blocks.  These
coefficients should be fitted empirically from no-trace fixture benchmarks.
They are the main way the model captures kernel launch overhead, cuBLAS
handle/stream effects, small-GEMM overhead, peer-copy setup overhead, and
duplicate first-stage group costs.

\subsection{Temporary Memory Pressure}

Let $M_d$ be the peak temporary memory required by first-stage temporaries,
temporary replicas, peer input replicas, and partial outputs on device $d$
under the chosen schedule.  The first optimizer can enforce a hard constraint
\begin{equation}
  M_d \le M^{\max}_d
\end{equation}
or add a soft penalty
\begin{equation}
  C^{\mathrm{memory}}_d = \mu M_d.
\end{equation}
The hard constraint is preferable for production scheduling.  The soft penalty
is useful while fitting a prototype model.

\section{Integer-Program View}

For a fixed exact temporary factorization, the restricted placement domain can
be linearized.  Introduce binary variables
\begin{align}
  o_{k,d} &= 1 \quad \text{if center block } k \text{ is owned by device } d,\\
  a_{i,d} &= 1 \quad \text{if environment block } A_i \text{ is resident on } d,\\
  c_{i,d} &= 1 \quad \text{if environment block } C_i \text{ is resident on } d,\\
  y^L_{pab,d} &= 1 \quad \text{if product } A_aB_b \text{ contributes to } P^L_p
                  \text{ on } d,\\
  y^R_{qbc,d} &= 1 \quad \text{if product } B_bC_c \text{ contributes to } P^R_q
                  \text{ on } d,\\
  s^L_{rpc,d} &= 1 \quad \text{if } P^L_pC_c \text{ for } R_r \text{ runs on } d,\\
  s^R_{raq,d} &= 1 \quad \text{if } A_aP^R_q \text{ for } R_r \text{ runs on } d.
\end{align}
The assignment constraints are
\begin{align}
  \sum_{d \in \mathcal{D}} o_{k,d} &= 1
    &&\text{for every } k \in \mathcal{K},\\
  \sum_{d \in \mathcal{D}} a_{i,d} &= 1
    &&\text{for every } i \in \mathcal{A},\\
  \sum_{d \in \mathcal{D}} c_{i,d} &= 1
    &&\text{for every } i \in \mathcal{C},\\
  \sum_{d \in \mathcal{D}} y^L_{pab,d} &= 1
    &&\text{for every nonzero } \lambda^L_{pab},\\
  \sum_{d \in \mathcal{D}} y^R_{qbc,d} &= 1
    &&\text{for every nonzero } \lambda^R_{qbc},\\
  \sum_{d \in \mathcal{D}} s^L_{rpc,d} &= 1
    &&\text{for every nonzero } \mu^L_{rpc},\\
  \sum_{d \in \mathcal{D}} s^R_{raq,d} &= 1
    &&\text{for every nonzero } \mu^R_{raq}.
\end{align}
Activation variables for $B$ replicas, $A/C$ staging, temporary migrations,
partial outputs, and peer reductions can be constrained above the corresponding
$y$ and $s$ variables.  The conservative specialization can add constraints
$s^L_{r,(r,c),c,d} = c_{c,d}$ and $s^R_{r,a,(r,a),d} = a_{a,d}$ so the final
temporary accumulation happens where the corresponding environment block is
resident.  The makespan objective is then
\begin{equation}
  \min T
  \qquad \text{subject to} \qquad
  C_d \le T \quad \text{for every } d \in \mathcal{D}.
\end{equation}
This form is a mixed-integer linear program once the cost coefficients are
fixed.  It can be used offline as an oracle for small fixtures.  The runtime
optimizer can use cheaper local search or graph partitioning heuristics against
the same evaluator.

\section{Relation to Current Implementation}

The current deterministic resident executor is a strict subcase:
\begin{align}
  \mathcal{P}_L &= \emptyset,\\
  q &= (b,c,d),\\
  P^R_q &= B_bC_c \quad \text{constructed separately for each device } d
            \text{ that needs it},\\
  \kappa_R(q,b,c) &= d,\\
  \sigma_R(r,a,q) &= h(r) = d.
\end{align}
That is, every term is right-first, each right-first temporary is local to the
owner of the output block that needs it, identical $(B_b,C_c)$ temporaries are
reused only within one device, and cross-device temporary migration is not
modelled.  The current \texttt{empirical-contiguous} policy is a legacy
two-device ordered-range diagnostic in the existing block construction order.
That ordering is not a mathematical invariant: block spaces can be permuted and
any chosen partition can be packed contiguously afterward.  Therefore an
ordered-range policy is not an appropriate optimization domain beyond
diagnostic comparison with the real hypergraph placement model.

The next model should separate logical block identity, placement, and storage
packing:
\begin{enumerate}
  \item choose an exact temporary factorization and placement variables over
        logical block and term keys;
  \item evaluate the induced copies, first-stage groups, partial outputs, and
        reductions;
  \item pack the chosen placement into coalesced device slabs as a storage
        implementation detail.
\end{enumerate}

\section{Environment Construction Variant}

For environment construction, the same term hypergraph and path labels apply,
but the Hamiltonian-specific constraint
\begin{equation}
  \operatorname{owner}(B_k) = \operatorname{owner}(R_k)
\end{equation}
is not required.  The more general domain uses separate owner maps
\begin{equation}
  h^B : \mathcal{B} \to \mathcal{D}, \qquad
  h^R : \mathcal{R} \to \mathcal{D}.
\end{equation}
The output may also be written to host-side cached environment storage after
construction rather than retained as the next Krylov input.  This changes the
output-reduction and staging terms, but not the underlying typed hypergraph.

\end{document}
